\section{Conclusions}
\label{sec:conclusion}

We applied Topological Data Analysis to characterize the grokking phenomenon in a small transformer trained on modular arithmetic tasks. Across two datasets, five training fractions, and four architectural layers, we found that successful generalization is strongly associated with a structural regularization of the learned representation manifold. Specifically, $\beta_0$ converges to a stable value in both the decoder and linear layers concurrently with the rise in validation accuracy; i.e. decreasing in the decoder as intermediate representations consolidate into fewer connected components, and increasing in the linear layer as output representations organize into well-separated clusters; while intrinsic dimension decreases progressively across all layers to values substantially lower than those observed in non-generalizing conditions. These patterns are robust across all ten experimental conditions, and the partial generalization observed in the multiplication task at $q = 15\%$ produces intermediate topological behavior, suggesting the transition is continuous rather than abrupt.

Several directions remain open. First, a systematic sensitivity analysis over the pipeline hyperparameters --- the neighborhood size $k$, the distance percentile threshold, and the persistence noise cutoff $\tau$ --- is needed to establish the robustness of the reported signatures under different construction choices. Second, the identified topological markers should be tested on algebraic extensions of the current setting, including other cyclic groups $\mathbb{Z}_n$, non-commutative structures such as permutation groups, and alternative modular operations. More broadly, since the TDA pipeline is task-agnostic, it could in principle be applied to other domains where delayed generalization has been observed, such as sparse parity learning or image classification; validating whether the same topological signatures accompany grokking beyond modular arithmetic remains a longer-term goal. Finally, the correlational nature of our findings motivates exploring topological regularization losses that directly encourage manifold simplification, which could serve both as a causal test and as a practical mechanism to accelerate or induce grokking. 